\documentclass[10pt,a4paper]{amsart}
\usepackage[margin=19mm]{geometry}
\usepackage{amsmath,amssymb,mathtools}
\usepackage[scr=rsfs,cal=euler]{mathalfa}
\usepackage{xfrac}
\usepackage[unicode,hidelinks,bookmarks=false]{hyperref}
\newcommand{\ie}{i.e.\ }
\newcommand{\cf}{cf.\ }
\newcommand{\resp}{respectively\ }
\newcommand{\Subsection}{\subsection}
\providecommand{\nobreakdash}{\nobreak\hbox{-}}
\setlength{\emergencystretch}{2em}
\multlinegap0pt
\usepackage{bm}
\usepackage{bbm}
\usepackage{dsfont}
\usepackage{xspace}
\usepackage{cancel}
\usepackage{mathtools}
\usepackage{paralist}
\usepackage{upgreek}
\usepackage{amsthm}
\makeatletter
\@ifundefined{thm}{\newtheorem{thm}{Thm}}{}
\@ifundefined{lem}{\newtheorem{lem}[thm]{Lemma}}{}
\makeatother
\newcommand{\CE}{\mathscr{E}}
\newcommand{\R}{\mathbb{R}}
\newcommand{\bc}{{ \mathcal{H} }}
\newcommand{\e}{\mathrm{e}}
\newcommand{\fe}{\mathfrak{E}}
\renewcommand{\H}{\mathbb{H}}
\renewcommand{\Re}{\mathop{\mathrm{Re}}}
\renewcommand{\d}{\mathrm{d}}
\renewcommand{\i}{\mathrm{i}}
\title[The Spherical Maximal Operators on Hyperbolic Spaces]{The Spherical Maximal Operators on Hyperbolic Spaces}
\author{Peng Chen, Minxing Shen, Yunxiang Wang, Lixin Yan}
\date{Source: arXiv:2408.02180v3 (CC BY 4.0)}
\begin{document}
\maketitle

\noindent\emph{Source and licence.} The Spherical Maximal Operators on Hyperbolic Spaces. Version arXiv:2408.02180v3 (CC BY 4.0), distributed under the Creative Commons Attribution 4.0 International licence, as recorded in the arXiv licence metadata for this exact version and shown on the abstract page https://arxiv.org/abs/2408.02180v3. Full rights notice: licenses/arxiv-2408.02180-CC-BY-4.0.txt.\par\smallskip

\noindent\textit{Editorial scope.} Counted unit: the Lem whose source label is \texttt{lem:5.3} in the article (source0.tex lines 926--931, source label \texttt{lem:5.3}), delivered together with the article's own complete original proof of it (source0.tex lines 933--982, one \texttt{proof} environment of 50 lines). No mathematics is written, repaired or re-derived: statement and proof are the article's own text line for line. Required context retained uncounted: lem:2.0 (lines 383--393); e2.2 (lines 397--399); lem:2.1 (lines 402--409); lem:2.3 (lines 519--549); e4.201 (lines 909--913). Every definition, display and label of those lines is retained unchanged. Omitted from this package and not redistributed: 1--382, 394--396, 400--401, 410--518, 550--908, 914--925, 932--932, 983--1518. Nothing inside a retained excerpt is omitted or altered: every retained body is delivered exactly as the article's lines are written, so no omission receipt of any kind accompanies this package. Citations: the retained excerpts carry no citation command at all, so no bibliography entry of the article is transcribed and no reference is redistributed. Reference adaptation: none was needed and none was made; every reference inside the delivered unit resolves inside it, so no unresolved reference or citation is produced. Printed slips kept literal: no printed word, symbol, hypothesis or number of the counted statement or of its proof was altered. Layout and class: the article's own class is \texttt{amsproc} with packages including amsthm, amssymb, paralist, mathrsfs, graphicx, tikz, bbm, upgreek, float, hyperref, ulem; the wrapper is the lane's pinned \texttt{amsart} editorial wrapper at 10pt on A4, which restates the theorem-like environments the retained text needs and adds the article's own macro definitions that the retained text needs (\texttt{\textbackslash CE}, \texttt{\textbackslash H}, \texttt{\textbackslash R}, \texttt{\textbackslash Re}, \texttt{\textbackslash bc}, \texttt{\textbackslash d}, \texttt{\textbackslash e}, \texttt{\textbackslash fe}, \texttt{\textbackslash i}), with extra wrapper packages (bm, bbm, dsfont, xspace, cancel, mathtools, paralist, upgreek). The delivered numbering is the unit's own. Assets: no retained span contains a figure environment, a graphics-inclusion command, a PDF-page inclusion, a table or a TikZ picture; no image file is copied into this package, the package declares no asset dependency and no image file is redistributed with it. Licence: version arXiv:2408.02180v3 of this article is distributed under the Creative Commons Attribution 4.0 International licence, as recorded in the arXiv licence metadata for this exact version and shown on the abstract page https://arxiv.org/abs/2408.02180v3. A full rights notice is delivered at licenses/arxiv-2408.02180-CC-BY-4.0.txt. Characters: every retained excerpt is pure ASCII, so the delivered engine, XeLaTeX with its default Latin Modern text font, reports no missing character.
\begin{lem}\label{lem:2.0}
Suppose $\lambda\in \R$.
The Harish-Chandra function $\bc(\lambda)$ satisfies
$
|\bc(\lambda)|^{-2}=(\bc(\lambda)\,\bc(-\lambda))^{-1}.
$
The function $ \lambda^{-1}(\bc(-\lambda))^{-1}$ belongs to $C^\infty(\R)$ and
\begin{eqnarray*}
\left|{\d^k\over\d\lambda^k}(\lambda^{-1}(\bc(\lambda))^{-1})\right|\leq C_k\,(1+|\lambda|)^{(n-1)/2-1-k}.
\end{eqnarray*}
\end{lem}

\begin{eqnarray}\label{e2.2}
\bc^\alpha(\lambda)={2^{n-2}\Gamma(n/2)\over \sqrt{\uppi}}{\Gamma(\i\lambda)\over\Gamma((n-1)/2+\alpha+\i\lambda)}, \ \ \  {\rm Re}\, \alpha>{(1-n)/2}.
\end{eqnarray} 

\begin{lem}\label{lem:2.1}
Let  ${\rm Re}\, \alpha>{(1-n)/2}$ and $k\in{\mathbb N}$.  Then the  function $  \lambda\,\bc^\alpha(\lambda)$ belongs to $C^\infty(\R)$ and
\begin{eqnarray*}
\left|{\d^k\over\d\lambda^k}(\lambda \,\bc^\alpha(\lambda))\right|\leq C_{\alpha,k}\,(1+|\lambda|)^{1-(n-1)/2-\Re\alpha-k}.
\end{eqnarray*}


\end{lem}

\begin{lem}\label{lem:2.3}
Suppose $\lambda\in\R$.
\begin{asparaenum}[(i)]
\item If $t>0$, then
\begin{eqnarray*}
|m^\alpha_t(\lambda)|\leq C_\alpha \,(1+t)\,\e^{-(n-1) t/2}.
\end{eqnarray*}

\item If $0<t\leq \uppi$ satisfying $t|\lambda|\geq 1$, then $m^\alpha_\lambda(t)$ can be written in the form
\begin{eqnarray}\label{emmm}
m^\alpha_t(\lambda)=\e^{\i\lambda t}a^{\alpha,N}_1(\lambda,t)+\e^{-\i\lambda t}a^{\alpha,N}_1(-\lambda,t)+E^{\alpha,N}(\lambda,t),
\end{eqnarray}
where
\begin{eqnarray}\label{e2.601}
\begin{cases}
\displaystyle|{\partial^k_\lambda}{\partial^l_t}a^{\alpha,N}_1(\lambda,t)|\leq C_N\,(|\lambda|t)^{-\Re \alpha-(n-1)/2}|\lambda|^{-k}t^{-l},\\[4pt]
|E^{\alpha,N}(\lambda,t)|\leq C_N\,(|\lambda|t)^{-\Re\alpha-(n-1)/2-N-1} 
\end{cases}
\end{eqnarray}
for all integers $N, k, l\geq 0$, and for $\lambda, t$ in the ranges stated above.

\item If $t>(\log 2)/2$, then $m^\alpha_t(\lambda)$ can be written in the form
\begin{eqnarray*}
m^\alpha_t(\lambda)=\e^{-(n-1) t/2}\left(\e^{\i\lambda t}\bc^\alpha(\lambda)\,a^\alpha_2(\lambda,t)+\e^{-\i\lambda t}\bc^\alpha(-\lambda)\,a^\alpha_2(-\lambda,t)\right),
\end{eqnarray*}
where $\bc^\alpha$ is as in \eqref{e2.2} and $a^\alpha_2(\lambda,t)\in S^0_\lambda$, i.e.,
\begin{eqnarray*}
|\partial^k_\lambda\partial^l_t a^\alpha_2(\lambda,t)|\leq C_{k,l}(1+|\lambda|)^{-k}.
\end{eqnarray*}
\end{asparaenum}
\end{lem}

\begin{eqnarray}\label{e4.201}
\CE^\alpha_{t,j}(r)&=&
\phi_0(r)\,\e^{-(n-1)t/2}  \notag\\
&&\times\int_0^\infty \beta_j(\lambda)\left( \e^{\i\lambda t}\bc^\alpha(\lambda)\,a^\alpha_2(\lambda,t)+\e^{-\i\lambda t}\bc^\alpha(-\lambda)\,a^\alpha_2(-\lambda,t)\right)\varphi_\lambda(r)\,|\bc(\lambda)|^{-2}\,\d \lambda.
\end{eqnarray}

\begin{lem}\label{lem:5.3}
For $1\leq p\leq\infty$, $0\leq \sigma<(n-1)/2$ and $j\geq 0$
\begin{eqnarray}\label{e00}
\|\fe^\alpha_{T,j}(f)\|_p\leq C_{\alpha,\sigma}2^{-(\Re\alpha+\sigma)j}\e^{-(n-1)T/2}\|f\|_p.
\end{eqnarray}
\end{lem}

\begin{proof}
By \eqref{e4.201} and Lemmas \ref{lem:2.1} and \ref{lem:2.3}(i)(iii),
\begin{eqnarray*}
\sup_{t\in[T,T+1]}|\CE^\alpha_{t,0}(r)|\leq C_\alpha\phi_0(r)\,\e^{-(n-1)T/2}.
\end{eqnarray*}
It follows from Schur's lemma  that \eqref{e00} holds  for the case $j=0$.

Now we assume $j\geq 1$. We take an even cutoff function $\gamma\in C^\infty_c(\R)$ with $\gamma\equiv 1$ on $[-1,1]$ and $\gamma\equiv0$ outside $(-2,2)$. By Lemma \ref{lem:2.3}(ii), we write
$$
\CE^\alpha_{t,j}(r)=\CE^{\alpha, 1}_{t,j}(r) +\CE^{\alpha, 2}_{t,j}(r)+ \CE^{\alpha, 3}_{t,j}(r)
$$
where
\begin{eqnarray*} 	
\CE^{\alpha, 1}_{t,j}(r) 
&=&\ \phi_0(r)\,\e^{-(n-1)t/2} \\
&&\times  \int_0^\infty \beta_j(\lambda)\,\gamma(\lambda r) \left( \e^{\i\lambda t}\bc^\alpha(\lambda)\,a^\alpha_2(\lambda,t)+\e^{-\i\lambda t}\bc^\alpha(-\lambda)\,a^\alpha_2(-\lambda,t)\right)\varphi_\lambda(r)\,|\bc(\lambda)|^{-2}\,\d \lambda,
\end{eqnarray*}
\begin{eqnarray*} 	
\CE^{\alpha, 2}_{t,j}(r) 
&=& \phi_0(r)\,\e^{-(n-1)t/2} \\ 
&&\times\int_\R\beta_j(\lambda)\,(1-\gamma(\lambda r))\left(\e^{\i\lambda (t+ r)}a^{0,N}_1(\lambda,r)+\e^{\i\lambda (t- r)}a^{0,N}_1(-\lambda,r)\right)\bc^\alpha(\lambda)\,a^\alpha_2(\lambda,t)\,|\bc(\lambda)|^{-2}\,\d \lambda,
\end{eqnarray*}

\begin{eqnarray*}	
\CE^{\alpha, 3}_{t,j}(r) 
&=&\phi_0(r)\,\e^{-(n-1)t/2}\\  
&&\times \int_0^\infty\beta_j(\lambda)\,(1-\gamma(\lambda r)) \left(\e^{\i\lambda t}\bc^\alpha(\lambda)\,a^\alpha_2(\lambda,t)+\e^{-\i\lambda t}\bc^\alpha(-\lambda)\,a^\alpha_2(-\lambda,t)\right)E^{\alpha,N}(\lambda,r)\,|\bc(\lambda)|^{-2}\,\d \lambda.
\end{eqnarray*} 	
By Lemmas \ref{lem:2.0}, \ref{lem:2.1} and \ref{lem:2.3} and integration by parts, taking $N$ sufficiently large, we obtain that $\CE^{\alpha, i}_{t,j}(r)\ (i=1,2,3)$ are all bounded by a constant multiple of $\phi_0(r)\,\e^{-(n-1)t/2}2^{-(\Re\alpha+\sigma)j}r^{-\sigma-(n+1)/2}$.
\iffalse
\begin{eqnarray*}
|\CE^{\alpha, 1}_{t,j}(r) | &\leq&  C_\alpha \phi_0(r)\,\e^{-(n-1)t/2}2^{-(\Re\alpha+\sigma)j}r^{-\sigma-(n+1)/2},\\
|\CE^{\alpha, 2}_{t,j}(r) |  &\leq&    C_{\alpha,N}\phi_0(r)\,\e^{-(n-1)t/2}\e^{-(n-1)t/2}2^{-(\Re\alpha+\sigma)j}r^{-\sigma-(n+1)/2+N},\\
|\CE^{\alpha, 3}_{t,j}(r) | &\leq&    C_{\alpha,N}\phi_0(r)\,\e^{-(n-1)t/2}\e^{-(n-1)t/2}2^{-(\Re\alpha+\sigma)j}r^{-\sigma-(n+1)/2}
\end{eqnarray*}
\fi
Hence, we obtain
\begin{eqnarray}\label{e000}
\sup_{t\in[T,T+1]}|\CE^\alpha_{t,j}(r)|\leq C_\alpha \phi_0(r)\,\e^{-(n-1)T/2}2^{-(\Re\alpha+\sigma)j}r^{-\sigma-(n+1)/2}.
\end{eqnarray}
Note that  $\Re\alpha>(1-n)/2$. By  using polar coordinates, one can derive
\begin{eqnarray*}
%&&\hspace{-1cm}
\sup_{z\in\H^n}\left\|\sup_{t\in[T,T+1]}|\CE^\alpha_{t,j}(d(z,\cdot))|\right\|_1+\sup_{w\in\H^n}\left\|\sup_{t\in[T,T+1]}|\CE^\alpha_{t,j}(d(\cdot,w))|\right\|_1
%\\&\leq &C_{\alpha}\e^{-(n-1)t/2}2^{-(\Re\alpha+\sigma) j}\int_0^2r^{-\sigma-(n+1)/2}(\sinh r)^{n-1}\,\d r
\leq C_{\alpha,\sigma}2^{-(\Re\alpha+\sigma)j}\e^{-(n-1)T/2}
\end{eqnarray*}
provided that  $\sigma<(n-1)/2$. Then we apply Schur's lemma to derive the desired result \eqref{e00}.
The proof of Lemma~\ref{lem:5.3} is concluded.
\end{proof}

\end{document}
